\begin{textsample}{POS Dim 1 – ro – Score 39.00 – ro/ro\_em\_4.txt}  \label{ex:f1_pos_150}
1.3.
Contexto do Ensino Médio O Ensino Médio, \textbf{enquanto} etapa conclusiva da Educação Básica, de acordo com artigo nº 35 da Lei de Diretrizes e Bases da Educação nº 9394/96, tem por objetivo atender os jovens oriundos do Ensino Fundamental e prepará -los para a conclusão do processo formativo no decorrer de \textbf{uma} duração mínima de três anos.
Para tanto, é \textbf{crucial} a organização de \textbf{uma} proposta pedagógica \textbf{que} proporcione ao estudante \textbf{uma} formação integral, ao compreender as determinações da vida social e produtiva ; articulando trabalho, ciência, tecnologia e cultura na perspectiva da emancipação humana.
Além disso, para consolidação e aprofundamento dos conhecimentos adquiridos na etapa do Ensino Fundamental, para fins de prosseguimento dos estudos e a preparação básica para o trabalho e a cidadania do estudante, é relevante também a compreensão dos \textbf{fundamentos} científico-tecnológicos dos processos produtivos, relacionando a \textbf{teoria} com a prática, no ensino de cada \textbf{disciplina}.
Para tanto, mostra -se como sendo de grande relevância o aprimoramento e desenvolvimento não só de competências cognitivas, mas também as competências socioemocionais \textbf{que} contribuam efetivamente para a formação ética e o desenvolvimento da autonomia intelectual e do pensamento crítico.
Com base nisso, cabe evidenciar \textbf{que} é na etapa do Ensino Médio \textbf{que} os jovens intensificam o conhecimento sobre seus sentimentos, interesses, capacidades intelectuais e \textbf{expressivas} ; ampliam e aprofundam vínculos sociais e afetivos ; e refletem sobre a vida, trabalho e a profissão \textbf{que} gostariam de ter e ser.
Encontram -se diante de questionamentos sobre si próprios e seus projetos de vida, \textbf{vivendo} juventudes marcadas por contextos culturais e sociais diversos.
Conforme aponta, Arroyo ( 2015 ), \textbf{que} “ é essencial o reconhecimento de \textbf{que} os jovens estudantes \textbf{que} vão chegando ao Ensino Médio são também Outros, de outras origens sociais, raciais, étnicas, dos campos e das periferias ”. ( ARROYO, 2014, p. 55 ).
São jovens com \textbf{uma} diversidade enorme, e se encontram todos em \textbf{um} único lugar.
Em vista disso, é importante evidenciar \textbf{que} o foco \textbf{central} da aprendizagem significativa não perpassa única e exclusivamente pelo interesse em formar o estudante para atender as exigências do mercado de trabalho, mas sim buscar \textbf{uma} formação sintonizada \textbf{que} venha prepará -lo para conquistar com \textbf{equidade} e de forma ética e democrática \textbf{uma} melhor qualidade de vida.
Nesse contexto, na perspectiva de formação integral, o estudante deve ser sujeito histórico do seu próprio ambiente, buscando desenvolver a consciência crítica \textbf{que} o leve a trilhar caminhos para construção de \textbf{um} mundo melhor.
Além das finalidades da educação nacional enunciadas no art. 205 na Constituição Federal e no Art. nº 2 da LDB, \textbf{que} têm como foco o pleno desenvolvimento da pessoa, a preparação para o exercício da cidadania e a qualificação para o trabalho, deve -se considerar integradamente o previsto nas metas 2 e 3 do PNE ( 2014-2024 ), \textbf{que} assegura, à criança de 6 a 14 anos e ao adolescente de 15 até 17 anos, concluírem a educação básica no tempo e idade certa, bem como todos os direitos fundamentais \textbf{inerentes} à pessoa, as oportunidades oferecidas para o desenvolvimento físico, mental, moral, espiritual e social, em condições de liberdade e de \textbf{dignidade}.
E a Emenda Constitucional nº 59, de 11 de novembro de 2009, \textbf{que} dá nova redação aos incisos I e VII do art. nº 208 da CF/88, de forma a prever a obrigatoriedade do ensino de quatro a dezessete anos.
As normas supracitadas asseguram a Educação Básica universal e o \textbf{alicerce} indispensável para desenvolver nas pessoas a capacidade de exercer em plenitude o direito à cidadania.
É o tempo, o espaço e o contexto em \textbf{que} o \textbf{sujeito} aprende a constituir e reconstituir a sua identidade, em meio a transformações corporais, afetivo emocionais, sócio emocionais, cognitivas e socioculturais, respeitando e valorizando as diferenças.
Da aquisição plena desse direito \textbf{depende} a possibilidade de exercitar todos os demais direitos definidos na Constituição, no Estatuto da Criança e do Adolescente ( ECA ), na legislação ordinária e nas \textbf{inúmeras} disposições legais \textbf{que} consagram as prerrogativas do cidadão brasileiro.
Somente \textbf{um} ser educado terá condição efetiva de participação social, ciente e consciente de seus direitos e deveres civis, sociais, políticos, econômicos e éticos.
Em consonância a essas diretrizes, o Art. 5º da Resolução CNE/CEB nº 3/2018 define \textbf{que} o ensino médio em todas as suas modalidades de ensino e as suas formas de organização e oferta, além dos princípios gerais estabelecidos para a educação nacional no art. 206 da Constituição Federal e no Art. nº 3º da LDB, será orientado pelos seguintes princípios específicos : I - formação integral do estudante, expressa por valores, aspectos físicos, cognitivos e socioemocionais ; II - projeto de vida como estratégia de reflexão sobre trajetória escolar na construção das dimensões pessoal, cidadã e profissional do estudante ; III - pesquisa como prática pedagógica para inovação, criação e construção de novos conhecimentos ; IV - respeito aos direitos humanos como direito universal ; V - compreensão da diversidade e realidade dos \textbf{sujeitos}, das formas de produção e de trabalho e das culturas ; VI - sustentabilidade ambiental ; VII - diversificação da oferta de forma a possibilitar múltiplas trajetórias por parte dos estudantes e a articulação dos saberes com o contexto histórico, econômico, social, científico, ambiental, cultural local e do mundo do trabalho ; VIII - indissociabilidade entre educação e prática social, considerando -se a historicidade dos conhecimentos e dos protagonistas do processo educativo ; IX - indissociabilidade entre \textbf{teoria} e prática no processo de \textbf{ensino-aprendizagem}. ( BRASIL, 2018 ).
Por \textbf{intermédio} desses princípios, busca -se estabelecer \textbf{um} diálogo com os jovens/estudantes ao direcionar o trabalho pedagógico a partir de ações em torno das mais diferentes expressões culturais, na perspectiva de valorizar a cultura juvenil dentro da comunidade escolar.
Para tanto, durante o processo de formação do estudante, é necessário estabelecer critérios voltados para a dimensão pessoal, social, cidadã e profissional do ser fundamentado em princípios éticos e morais.
Referente aos princípios \textbf{metodológicos} \textbf{que} orientam a proposta pedagógica para o Ensino Médio no Estado de Rondônia, é pertinente destacar \textbf{que} em conformidade com a Base Nacional Comum Curricular, os arranjos da Educação Básica na \textbf{contemporaneidade}, têm por finalidade oportunizar as condições para mobilização de saberes e a articulação de \textbf{uma} aprendizagem centrada na figura do estudante \textbf{que} contribuam para a significação dessa formação em sua vida pessoal e profissional.
E, assim, como princípios básicos, o ato de \textbf{educar} deve oportunizar mudanças de perspectivas na forma de conhecer, aprender e absorver o conhecimento e nas formas de ser e agir do estudante.
Partindo desses \textbf{fundamentos}, defende -se como ponto de partida, o aprofundamento \textbf{teórico} e prático em metodologias \textbf{que} \textbf{consigam} contribuir para a transformação dos saberes científicos em saberes escolares significativos para o cotidiano do estudante. \textbf{Ademais}, como pressupostos \textbf{teóricos} para integração curricular no Ensino Médio, a Educação por competências \textbf{revela} -se como sendo \textbf{uma} grande aliada quando se \textbf{fomenta} o ato de ensinar, e aprender a aprender a partir de situações cotidianas e/ou reais, pois o \textbf{que} é ensinado dentro da escola deve acompanhar todo esse movimento da sociedade.
Vale ressaltar \textbf{que} este novo formato \textbf{que} se \textbf{configura}, chega para romper com o ensino tradicional e oportunizar a construção de novas formas de aprender a aprender e aprender a ensinar, de modo \textbf{que}, é necessário reconhecer, refletir e ressignificar a aprendizagem em meio a esse turbilhão \textbf{que} \textbf{permeia} a vida na \textbf{contemporaneidade}.
Nesses termos, entende -se como proposta principal do ensino por competências “ facilitar a capacidade de transferir aprendizagens \textbf{que} geralmente foram apresentadas descontextualizadas, às situações próximas da realidade, o \textbf{que} representa \textbf{uma} redefinição do objeto de estudo da escola ” ( ZABALA ; ARNAU, 2010, p. 110 ).
Sendo assim, no \textbf{que} \textbf{diz} respeito à \textbf{aplicabilidade} da Educação por competências no Ensino Médio, o relatório da UNESCO para Comissão Internacional \textbf{que} trata sobre a Educação para o \textbf{século} XXI, ao evidenciar os impactos e as consequências da sociedade da informação e do conhecimento, constatou a necessidade de \textbf{uma} educação continuada sob a qual seja garantida ao estudante a aprendizagem ao longo da vida considerando os quatros pilares da Educação, sendo eles : aprender a conhecer ; aprender a fazer ; aprender a \textbf{viver} juntos e aprender a ser ( MORAN ; MASETTO ; BEHRENS, 2013 ).
Em vista disso, a proposta pedagógica para o ensino médio do Estado de Rondônia aponta o campo das metodologias ativas com sendo \textbf{um} caminho \textbf{metodológico} e estratégico relevante e fundamental para reflexão na ação e sobre a ação ( SCHÖN, 2000 ) no \textbf{que} \textbf{diz} respeito a re(criar) experiências educativas significativas para o estudante do Ensino Médio nas áreas, entre as áreas de conhecimentos e na parte flexível do currículo. \textbf{Ademais}, defende a importância da construção de \textbf{uma} prática educativa fundamentada na Educação por competência \textbf{que} forneça bases para \textbf{que} o professor possa compreender as mudanças nas concepções do aprender, planejar, ensinar e comunicar na sociedade da informação, do conhecimento e da aprendizagem.
Atrelados a isso, a produção do saber nas áreas do conhecimento demanda pelo desenvolvimento de ações \textbf{que} levem o professor e o estudante a buscar os processos de investigação e pesquisa com vistas a abrir caminhos coletivos para produção de conhecimento de forma integrada por parte do aluno e professor.
Nesse âmbito, torna -se importante \textbf{que} o estudante se desvincule de ações passivas como escutar, ler, decorar e de ser repetidor fiel dos ensinamentos do professor, para se tornar criativo, crítico, \textbf{pesquisador} e ativo na produção do conhecimento. \textbf{Ademais}, o trabalho pedagógico com temas \textbf{contemporâneos} é \textbf{crucial} para \textbf{explicitar} a conexão entre os diferentes componentes curriculares das áreas de conhecimento.
Nesse sentido, o conhecimento deve ser apresentado de forma integrada, articulada e contextualizada com as situações vivenciadas pelos estudantes em suas realidades, mediante a \textbf{uma} \textbf{abordagem} transversal e interdisciplinar dos objetos do conhecimento descritos na BNCC.
Dessa forma, o Ensino Médio, como etapa conclusiva da Educação Básica, descrita na LDB 9394/96, no artigo 35, tem por objetivo atender os jovens oriundos do ensino fundamental como preparação para a conclusão do processo formativo, com a duração mínima de três anos. \textbf{Posto} isso, o ensino médio, é \textbf{uma} etapa responsável pela consolidação e conclusão do processo formativo da Educação Básica e deve se organizar para proporcionar ao estudante \textbf{uma} formação integral, conduzindo -o a pensar e compreender as determinações da vida social e produtiva.
Nessa perspectiva, o currículo do ensino médio, deve ter por \textbf{escopo} a redução da distância entre as atividades escolares e as práticas sociais, apresentando \textbf{uma} base unitária sobre a qual podem se assentar possibilidades diversas, como preparação geral para prosseguimento dos estudos ou, facultativamente, preparação para profissões técnicas.
Por exemplo, na ciência e na tecnologia, como iniciação científica e tecnológica ; nas artes e na cultura, como ampliação da formação cultural, possibilitando a flexibilização necessária \textbf{que} permita ao estudante optar por áreas de conhecimento, visando a continuidade de estudos ou optar pela formação técnica e profissional, objetivando inserir -se imediatamente, após a conclusão, no mundo do trabalho.
Apesar de todo esforço empreendido nesse sentido, o conjunto de funções atribuídas ao Ensino Médio não tem atendido às necessidades dos jovens dos dias \textbf{atuais} e nem atende às demandas do mercado de trabalho.
Portanto, para \textbf{que} se assegure a permanência dos jovens na escola, com proveito, até a conclusão da Educação Básica, os sistemas educativos devem prever currículos flexíveis, com diferentes alternativas, para \textbf{que} os jovens tenham a oportunidade de escolher o percurso formativo \textbf{que} mais atenda aos seus interesses, suas necessidades e suas aspirações.
E, segundo o documento da UNICEF “ 10 Desafios do Ensino Médio no Brasil para garantir o direito de aprender de adolescentes de 15 A 17 anos ” ( BRASÍLIA, 2014 ), \textbf{um} dos grandes desafios para o Brasil no \textbf{que} \textbf{diz} respeito à garantia dos direitos de seus adolescentes é a universalização do Ensino médio – etapa adequada para a faixa etária de 15 a 17 anos, \textbf{que} se tornou obrigatória a partir da Emenda Constitucional nº 59, de 2009.
Independentemente do lugar, a relação dos adolescentes com a escola é muito parecida e os obstáculos também são semelhantes.
Alguns deles estão relacionados com o contexto socioeconômico, como o trabalho precoce, a gravidez e a violência familiar.
Outros estão vinculados às questões relacionadas à organização da escola, como os conteúdos distantes da realidade dos alunos, a falta de diálogo entre alunos, professores e a gestão da escola, a desmotivação, a violência do cotidiano escolar e a infraestrutura precária dos estabelecimentos.

% matched lemmas: abordagem, ademais, alicerce, aplicabilidade, atual, central, configurar, conseguir, contemporaneidade, contemporâneo, crucial, depender, dignidade, disciplina, dizer, educar, enquanto, ensino-aprendizagem, equidade, escopo, explicitar, expressivo, fomentar, fundamento, inerente, intermédio, inúmero, metodológico, permear, pesquisador, postar, que, revelar, sujeito, século, teoria, teórico, um, viver
\end{textsample}
